\documentclass{amsart}
\usepackage[UTF8,fontset=windows]{ctex}
\usepackage[a4paper,margin=24mm]{geometry}
% For dealing with references we use the comment environment
\usepackage{verbatim}
\newenvironment{reference}{\comment}{\endcomment}
%\newenvironment{reference}{}{}
\newenvironment{slogan}{\comment}{\endcomment}
\newenvironment{history}{\comment}{\endcomment}

% For commutative diagrams we use Xy-pic
\usepackage[all]{xy}

% We use 2cell for 2-commutative diagrams.
\xyoption{2cell}
\UseAllTwocells

% We use multicol for the list of chapters between chapters
\usepackage{multicol}

% This is generally recommended for better output
\usepackage{lmodern}
\usepackage[T1]{fontenc}

% For cross-file-references

% Package for hypertext links:
\usepackage{hyperref}

% For any local file, say "hello.tex" you want to link to please

% Theorem environments.
%
\theoremstyle{plain}
\newtheorem{theorem}[subsection]{Theorem}
\newtheorem{proposition}[subsection]{Proposition}
\newtheorem{lemma}[subsection]{Lemma}

\theoremstyle{definition}
\newtheorem{definition}[subsection]{Definition}
\newtheorem{example}[subsection]{Example}
\newtheorem{exercise}[subsection]{Exercise}
\newtheorem{situation}[subsection]{Situation}

\theoremstyle{remark}
\newtheorem{remark}[subsection]{Remark}
\newtheorem{remarks}[subsection]{Remarks}

\numberwithin{equation}{subsection}

% Macros
%
\def\lim{\mathop{\mathrm{lim}}\nolimits}
\def\colim{\mathop{\mathrm{colim}}\nolimits}
\def\Spec{\mathop{\mathrm{Spec}}}
\def\Hom{\mathop{\mathrm{Hom}}\nolimits}
\def\Ext{\mathop{\mathrm{Ext}}\nolimits}
\def\SheafHom{\mathop{\mathcal{H}\!\mathit{om}}\nolimits}
\def\SheafExt{\mathop{\mathcal{E}\!\mathit{xt}}\nolimits}
\def\Sch{\mathit{Sch}}
\def\Mor{\mathop{\mathrm{Mor}}\nolimits}
\def\Ob{\mathop{\mathrm{Ob}}\nolimits}
\def\Sh{\mathop{\mathit{Sh}}\nolimits}
\def\NL{\mathop{N\!L}\nolimits}
\def\CH{\mathop{\mathrm{CH}}\nolimits}
\def\proetale{{pro\text{-}\acute{e}tale}}
\def\etale{{\acute{e}tale}}
\def\QCoh{\mathit{QCoh}}
\def\Ker{\mathop{\mathrm{Ker}}}
\def\Im{\mathop{\mathrm{Im}}}
\def\Coker{\mathop{\mathrm{Coker}}}
\def\Coim{\mathop{\mathrm{Coim}}}

% Boxtimes
%
\DeclareMathSymbol{\boxtimes}{\mathbin}{AMSa}{"02}

%
% Macros for moduli stacks/spaces
%
\def\QCohstack{\mathcal{QC}\!\mathit{oh}}
\def\Cohstack{\mathcal{C}\!\mathit{oh}}
\def\Spacesstack{\mathcal{S}\!\mathit{paces}}
\def\Quotfunctor{\mathrm{Quot}}
\def\Hilbfunctor{\mathrm{Hilb}}
\def\Curvesstack{\mathcal{C}\!\mathit{urves}}
\def\Polarizedstack{\mathcal{P}\!\mathit{olarized}}
\def\Complexesstack{\mathcal{C}\!\mathit{omplexes}}
% \Pic is the operator that assigns to X its picard group, usage \Pic(X)
% \Picardstack_{X/B} denotes the Picard stack of X over B
% \Picardfunctor_{X/B} denotes the Picard functor of X over B
\def\Pic{\mathop{\mathrm{Pic}}\nolimits}
\def\Picardstack{\mathcal{P}\!\mathit{ic}}
\def\Picardfunctor{\mathrm{Pic}}
\def\Deformationcategory{\mathcal{D}\!\mathit{ef}}

\hypersetup{hidelinks}
\widowpenalty=10000\clubpenalty=10000\displaywidowpenalty=10000
\begin{document}
\title[Stacks Project, Tag 0GI9]{形式态射 versality 的统一光滑邻域}
\maketitle
\noindent\textbf{Editorial reading note.} The printed reversal A\_n to B\_n in the final sentence is retained. The preceding argument concerns B\_n to A\_n.


\noindent\textbf{Counting scope.} The primary problem is Tag 0GI9, marked below. All preceding numbered statements are supporting author context and are not counted as additional problems. This is a local pilot with a temporary ID, not production admission.


\noindent Copyright (C) 2005--2025 Johan de Jong. Stacks Project authors. Source: \href{https://stacks.math.columbia.edu/tag/0GI9}{Tag 0GI9}.

\noindent Difficulty H2: 本题从一点versality推出存在一个对所有形式厚化层同时有效的光滑开邻域。各层分别光滑只给随n变化的开集；证明把所有分次层装入Rees代数映射，并利用核余核有限性在一次局部化后同时成为同构。这个统一化步骤及étale代数化归约仍由本证明完成，判H2。.

\noindent History: excerpt from revision local cached snapshot; exact chapter hashes; revision unknown, artin.tex, lines 6112--6232. Only reference presentation was adapted. Licensed under GNU FDL 1.2 or later; no invariant sections or cover texts. See accompanying COPYING.

\section{Axioms for functors}
\label{section-axioms-functors}

\noindent
Let $S$ be a scheme. Let $F : (\Sch/S)_{fppf}^{opp} \to \textit{Sets}$ be a
functor. Denote $\mathcal{X} = \mathcal{S}_F$ the category fibred in sets
associated to $F$, see Algebraic Stacks, Section \href{https://stacks.math.columbia.edu/tag/04SU}{section split (Tag 04SU)}.
In this section we provide a translation between the material above
as it applies to $\mathcal{X}$, to statements about $F$.

\medskip\noindent
Let $S$ be a locally Noetherian scheme. Let
$F : (\Sch/S)_{fppf}^{opp} \to \textit{Sets}$ be a functor. Let $k$ be
a field of finite type over $S$. Let $x_0 \in F(\Spec(k))$.
The associated predeformation category (\href{https://stacks.math.columbia.edu/tag/07T4}{equation predeformation category (Tag 07T4)})
corresponds to the functor
$$
F_{k, x_0} : \mathcal{C}_\Lambda \longrightarrow \textit{Sets},
\quad
A \longmapsto \{x \in F(\Spec(A)) \mid x|_{\Spec(k)} = x_0 \}.
$$
Recall that we do not distinguish between
categories cofibred in sets over $\mathcal{C}_\Lambda$
and functor $\mathcal{C}_\Lambda \to \textit{Sets}$,
see Formal Deformation Theory, Remarks
\href{https://stacks.math.columbia.edu/tag/06GK}{remarks cofibered groupoids (Tag 06GK)}
(\href{https://stacks.math.columbia.edu/tag/06GP}{item convention cofibered sets (Tag 06GP)}).
Given a transformation of functors $a : F \to G$, setting
$y_0 = a(x_0)$ we obtain a morphism
$$
F_{k, x_0} \longrightarrow G_{k, y_0}
$$
see (\href{https://stacks.math.columbia.edu/tag/07WJ}{equation functoriality (Tag 07WJ)}).
Lemma \href{https://stacks.math.columbia.edu/tag/07WK}{lemma formally smooth on deformation categories (Tag 07WK)} tells us that if
$a : F \to G$ is formally smooth (in the sense of
More on Morphisms of Spaces, Definition
\ref{spaces-more-morphisms-definition-formally-smooth-etale-unramified}), then
$F_{k, x_0} \longrightarrow G_{k, y_0}$ is smooth as
in Formal Deformation Theory, Remark
\href{https://stacks.math.columbia.edu/tag/06HJ}{remark compare smooth schlessinger (Tag 06HJ)}.

\medskip\noindent
Lemma \href{https://stacks.math.columbia.edu/tag/07WN}{lemma pushout (Tag 07WN)} says that if $Y' = Y \amalg_X X'$ in the
category of schemes over $S$ where $X \to X'$ is a thickening and
$X \to Y$ is affine, then the map
$$
F(Y \amalg_X X') \to F(Y) \times_{F(X)} F(X')
$$
is a bijection, provided that $F$ is an algebraic space.
We say a general functor $F$ satisfies the {\it Rim-Schlessinger condition}
or we say $F$ {\it satisfies (RS)} if given any
pushout $Y' = Y \amalg_X X'$ where $Y, X, X'$ are spectra of Artinian
local rings of finite type over $S$, then
$$
F(Y \amalg_X X') \to F(Y) \times_{F(X)} F(X')
$$
is a bijection. Thus every algebraic space satisfies (RS).

\medskip\noindent
Lemma \href{https://stacks.math.columbia.edu/tag/07WU}{lemma deformation category (Tag 07WU)} says that
given a functor $F$ which satisfies (RS), then all $F_{k, x_0}$
are deformation functors as in
Formal Deformation Theory, Definition
\ref{formal-defos-definition-deformation-category}, i.e., they satisfy
(RS) as in
Formal Deformation Theory, Remark
\href{https://stacks.math.columbia.edu/tag/06J6}{remark compare schlessinger H4 (Tag 06J6)}.
In particular the tangent space
$$
TF_{k, x_0} = \{x \in F(\Spec(k[\epsilon])) \mid x|_{\Spec(k)} = x_0\}
$$
has the structure of a $k$-vector space by Formal Deformation Theory,
Lemma \href{https://stacks.math.columbia.edu/tag/06IH}{lemma tangent space vector space (Tag 06IH)}.

\medskip\noindent
Lemma \href{https://stacks.math.columbia.edu/tag/07X1}{lemma finite dimension (Tag 07X1)} says that an algebraic space $F$
locally of finite type over $S$ gives rise to deformation functors
$F_{k, x_0}$ with finite dimensional tangent spaces $TF_{k, x_0}$.

\medskip\noindent
A {\it formal object}\footnote{This is what Artin calls a formal deformation.}
$\xi = (R, \xi_n)$ of $F$ consists of a Noetherian
complete local $S$-algebra $R$ whose residue field is of finite type
over $S$, together with elements $\xi_n \in F(\Spec(R/\mathfrak m^n))$
such that $\xi_{n + 1}|_{\Spec(R/\mathfrak m^n)} = \xi_n$. A formal
object $\xi$ defines a formal object $\xi$ of $F_{R/\mathfrak m, \xi_1}$.
We say $\xi$ is {\it versal} if and only if it is versal in the sense of
Formal Deformation Theory, Definition \ref{formal-defos-definition-versal}.
A formal object $\xi = (R, \xi_n)$ is called {\it effective}
if there exists an $x \in F(\Spec(R))$ such that
$\xi_n = x|_{\Spec(R/\mathfrak m^n)}$ for all $n \geq 1$.
Lemma \href{https://stacks.math.columbia.edu/tag/07X8}{lemma effective (Tag 07X8)} says that if $F$ is an algebraic space,
then every formal object is effective.

\medskip\noindent
Let $U$ be a scheme locally of finite type over $S$ and let $x \in F(U)$.
Let $u_0 \in U$ be a finite type point. We say that $x$ is versal at $u_0$
if and only if
$\xi = (\mathcal{O}_{U, u_0}^\wedge,
x|_{\Spec(\mathcal{O}_{U, u_0}/\mathfrak m_{u_0}^n)})$
is a versal formal object in the sense described above.

\medskip\noindent
Let $S$ be a locally Noetherian scheme. Let
$F : (\Sch/S)_{fppf}^{opp} \to \textit{Sets}$ be a functor.
Here are the axioms we will consider on $F$.
\begin{enumerate}
\item[{[-1]}]  a set theoretic condition\footnote{The condition is the
following: the supremum of all the cardinalities
$|F(\Spec(k))|$ where $k$ runs over the finite
type fields over $S$ is $\leq$ than the size of some
object of $(\Sch/S)_{fppf}$.} to be ignored by
readers who are not interested in set theoretical issues,
\item[{[0]}] $F$ is a sheaf for the \'etale topology,
\item[{[1]}] $F$ is limit preserving,
\item[{[2]}] $F$ satisfies the Rim-Schlessinger condition (RS),
\item[{[3]}] every tangent space $TF_{k, x_0}$ is finite dimensional,
\item[{[4]}] every formal object is effective,
\item[{[5]}] $F$ satisfies openness of versality.
\end{enumerate}
Here {\it limit preserving} is the notion defined in
Limits of Spaces, Definition
\ref{spaces-limits-definition-locally-finite-presentation} and
{\it openness of versality} means the following: Given a scheme $U$
locally of finite type over $S$, given $x \in F(U)$, and given
a finite type point $u_0 \in U$ such that $x$ is versal at $u_0$,
then there exists an open neighbourhood $u_0 \in U' \subset U$
such that $x$ is versal at every finite type point of $U'$.








\begin{lemma}
\label{lemma-base-change-versal}
Let $S$, $\mathcal{X}$, $U$, $x$, $u_0$ be as in
Definition \ref{definition-versal}. Assume
\begin{enumerate}
\item $\Delta : \mathcal{X} \to \mathcal{X} \times \mathcal{X}$
is representable by algebraic spaces,
\item $\Delta$ is locally of finite type
(for example if $\mathcal{X}$ is limit preserving), and
\item $\mathcal{X}$ has (RS).
\end{enumerate}
Let $V$ be a scheme locally of finite type over $S$
and let $y$ be an object of $\mathcal{X}$ over $V$.
Form the $2$-fibre product
$$
\xymatrix{
\mathcal{Z} \ar[r] \ar[d] & (\Sch/U)_{fppf} \ar[d]^x \\
(\Sch/V)_{fppf} \ar[r]^y & \mathcal{X}
}
$$
Let $Z$ be the algebraic space representing $\mathcal{Z}$
and let $z_0 \in |Z|$ be a finite type point lying over $u_0$.
If $x$ is versal at $u_0$, then
the morphism $Z \to V$ is smooth at $z_0$.
\end{lemma}

\begin{lemma}
\label{restricted-lemma-algebraize-rig-etale-affine}
Let $T \subset X$ be a closed subset of a Noetherian affine scheme $X$.
Let $W$ be a Noetherian affine formal algebraic space.
Let $g : W \to X_{/T}$ be a rig-\'etale morphism. Then there exists
an affine scheme $X'$ and a finite type morphism $f : X' \to X$
\'etale over $X \setminus T$ such that there is an isomorphism
$X'_{/f^{-1}T} \cong W$ compatible with $f_{/T}$ and $g$.
Moreover, if $W \to X_{/T}$ is \'etale, then $X' \to X$ is \'etale.
\end{lemma}

\begin{proof}
The existence of $X'$ is a restatement of
Lemma \href{https://stacks.math.columbia.edu/tag/0AKG}{lemma approximate by etale over complement (Tag 0AKG)}.
The final statement follows from
More on Morphisms, Lemma
\href{https://stacks.math.columbia.edu/tag/0A43}{lemma check smoothness on infinitesimal nbhds (Tag 0A43)}.
\end{proof}


\begin{definition}
\label{spaces-more-morphisms-definition-formally-smooth-etale-unramified}
Let $S$ be a scheme.
Let $a : F \to G$ be a transformation of functors
$F, G : (\Sch/S)_{fppf}^{opp} \to \textit{Sets}$.
Consider commutative solid diagrams of the form
$$
\xymatrix{
F \ar[d]_a & T \ar[d]^i \ar[l] \\
G & T' \ar[l] \ar@{-->}[lu]
}
$$
where $T$ and $T'$ are affine schemes and $i$ is a closed immersion
defined by an ideal of square zero.
\begin{enumerate}
\item We say $a$ is {\it formally smooth} if given any solid
diagram as above there exists a dotted arrow making the diagram
commute\footnote{This is just one possible definition that one can
make here. Another slightly weaker condition would be to require that
the dotted arrow exists fppf locally on $T'$. This weaker notion
has in some sense better formal properties.}.
\item We say $a$ is {\it formally \'etale} if given any solid
diagram as above there exists exactly one dotted arrow making the diagram
commute.
\item We say $a$ is {\it formally unramified} if given any solid
diagram as above there exists at most one dotted arrow making the diagram
commute.
\end{enumerate}
\end{definition}

\begin{definition}
\label{formal-defos-definition-deformation-category}
A {\it deformation category} is a predeformation category $\mathcal{F}$
satisfying (RS). A morphism of deformation categories is a morphism of
categories over $\mathcal{C}_\Lambda$.
\end{definition}

\begin{definition}
\label{formal-defos-definition-versal}
Let $\mathcal{F}$ be a category cofibered in groupoids.  Let $\xi$ be a formal
object of $\mathcal{F}$ lying over $R \in \Ob(\widehat{\mathcal{C}}_\Lambda)$.
We say $\xi$ is {\it versal} if the corresponding morphism
$\underline{\xi}: \underline{R}|_{\mathcal{C}_\Lambda} \to \mathcal{F}$
of Remark \href{https://stacks.math.columbia.edu/tag/06HC}{remark formal objects yoneda (Tag 06HC)} is smooth.
\end{definition}

\begin{definition}
\label{spaces-limits-definition-locally-finite-presentation}
Let $S$ be a scheme.
\begin{enumerate}
\item A functor $F : (\Sch/S)_{fppf}^{opp} \to \textit{Sets}$
is said to be {\it limit preserving} or {\it locally of finite presentation} if
for every affine scheme $T$ over $S$ which is a limit $T = \lim T_i$
of a directed inverse system of affine schemes $T_i$ over $S$, we have
$$
F(T) = \colim F(T_i).
$$
We sometimes say that $F$ is {\it locally of finite presentation over $S$}.
\item Let $F, G : (\Sch/S)_{fppf}^{opp} \to \textit{Sets}$.
A transformation of functors $a : F \to G$
is {\it limit preserving} or {\it locally of finite presentation}
if for every scheme $T$ over $S$ and every $y \in G(T)$ the functor
$$
F_y : (\Sch/T)_{fppf}^{opp} \longrightarrow \textit{Sets}, \quad
T'/T \longmapsto \{x \in F(T') \mid a(x) = y|_{T'}\}
$$
is locally of finite presentation over $T$\footnote{The characterization (2) in
Lemma \href{https://stacks.math.columbia.edu/tag/06BC}{lemma characterize relative limit preserving (Tag 06BC)}
may be easier to parse.}. We sometimes say that
$F$ is {\it relatively limit preserving} over $G$.
\end{enumerate}
\end{definition}

\begin{definition}
\label{definition-versal}
Let $S$ be a locally Noetherian scheme.
Let $\mathcal{X}$ be fibred in groupoids over $(\Sch/S)_{fppf}$.
Let $U$ be a scheme locally of finite type over $S$.
Let $x$ be an object of $\mathcal{X}$ lying over $U$.
Let $u_0$ be a finite type point of $U$.
We say $x$ is {\it versal} at $u_0$ if the morphism $\hat x$
(\href{https://stacks.math.columbia.edu/tag/07XE}{equation hat x (Tag 07XE)}) is smooth, see Formal Deformation Theory, Definition
\ref{formal-defos-definition-smooth-morphism}.
\end{definition}

\begin{definition}
\label{bootstrap-definition-property-transformation}
Let $S$ be a scheme. Let $a : F \to G$ be a map of presheaves on
$(\Sch/S)_{fppf}$ which is representable by algebraic spaces.
Let $\mathcal{P}$ be a property of morphisms of algebraic spaces which
\begin{enumerate}
\item is preserved under any base change, and
\item is fppf local on the base, see
Descent on Spaces,
Definition \ref{spaces-descent-definition-property-morphisms-local}.
\end{enumerate}
In this case we say that $a$ has {\it property $\mathcal{P}$} if for every
scheme $U$ and $\xi : U \to G$ the resulting morphism of algebraic spaces
$U \times_G F \to U$ has property $\mathcal{P}$.
\end{definition}

\begin{definition}
\label{formal-defos-definition-smooth-morphism}
Let $\varphi : \mathcal{F} \to \mathcal{G}$ be a morphism of categories
cofibered in groupoids over $\mathcal{C}_\Lambda$.  We say  $\varphi$ is
{\it smooth} if it satisfies the following condition: Let $B \to A$ be
a surjective ring map in $\mathcal{C}_\Lambda$.  Let $y \in
\Ob(\mathcal{G}(B)), x \in \Ob(\mathcal{F}(A))$, and $y
\to \varphi(x)$ be a morphism lying over $B \to A$.  Then there
exists $x' \in \Ob(\mathcal{F}(B))$, a morphism $x' \to x$
lying over $B \to A$, and a morphism $\varphi(x') \to y$ lying
over $\text{id}: B \to B$, such that the diagram
$$
\xymatrix{
\varphi(x') \ar[r] \ar[dr] & y \ar[d] \\
& \varphi(x)
}
$$
commutes.
\end{definition}

\begin{definition}
\label{spaces-descent-definition-property-morphisms-local}
Let $S$ be a scheme.
Let $\mathcal{P}$ be a property of morphisms of algebraic spaces over $S$.
Let $\tau \in \{fpqc, fppf, syntomic, smooth, \etale\}$.
We say $\mathcal{P}$ is {\it $\tau$ local on the base}, or
{\it $\tau$ local on the target}, or
{\it local on the base for the $\tau$-topology} if for any
$\tau$-covering $\{Y_i \to Y\}_{i \in I}$ of algebraic spaces
and any morphism of algebraic spaces $f : X \to Y$ we
have
$$
f \text{ has }\mathcal{P}
\Leftrightarrow
\text{each }Y_i \times_Y X \to Y_i\text{ has }\mathcal{P}.
$$
\end{definition}

\section*{Primary problem: Tag 0GI9}
\begin{lemma}
\label{lemma-openness-smoothness}
Let $S$ be a locally Noetherian scheme. Let $V$ be a scheme locally
of finite type over $S$. Let $Z \subset V$ be closed. Let $W$ be
a locally Noetherian formal algebraic space over $S$ such that
$W_{red}$ is locally of finite type over $S$. Let $g : V_{/Z} \to W$
be an adic morphism of formal algebraic spaces over $S$. Let $v \in V$
be a closed point such that $g$ is versal at $v$ (as in
Section \ref{section-axioms-functors}).
Then after replacing $V$ by an open neighbourhood of $v$ the
morphism $g$ is smooth (see proof).
\end{lemma}

\begin{proof}
Since $g$ is adic it is representable by algebraic spaces (Formal Spaces,
Section \href{https://stacks.math.columbia.edu/tag/0AQ2}{section adic (Tag 0AQ2)}).
Thus by saying $g$ is smooth we mean that $g$ should be smooth
in the sense of
Bootstrap, Definition \ref{bootstrap-definition-property-transformation}.

\medskip\noindent
Write $W = \colim W_n$ as in Formal Spaces, Lemma
\href{https://stacks.math.columbia.edu/tag/0GHM}{lemma structure locally noetherian (Tag 0GHM)}.
Set $V_n = V_{/Z} \times_{\hat x, W} W_n$.
Then $V_n$ is a closed subscheme with underlying set $Z$.
Smoothness of $V \to W$ is equivalent to the smoothness
of all the morphisms $V_n \to W_n$ (this holds because any morphism
$T \to W$ with $T$ a quasi-compact scheme factors through
$W_n$ for some $n$). We know that the morphism $V_n \to W_n$
is smooth at $v$ by
Lemma \ref{lemma-base-change-versal}\footnote{The lemma applies since
the diagonal of $W$ is representable by algebraic spaces and
locally of finite type, see Formal Spaces, Lemma
\href{https://stacks.math.columbia.edu/tag/0AN2}{lemma diagonal morphism formal algebraic spaces (Tag 0AN2)}
and we have seen that $W$ has (RS) in the proof of Lemma \href{https://stacks.math.columbia.edu/tag/0GI6}{lemma rs (Tag 0GI6)}.}.
Of course this means that given any $n$ we can shrink $V$
such that $V_n \to W_n$ is smooth. The problem is to find
an open which works for all $n$ at the same time.

\medskip\noindent
The question is local on $V$, hence we may assume $S = \Spec(R)$ and
$V = \Spec(A)$ are affine.

\medskip\noindent
In this paragraph we reduce to the case where $W$ is an affine formal
algebraic space. Choose an affine formal scheme $W'$ and an \'etale morphism
$W' \to W$ such that the image of $v$ in $W_{red}$ is in the
image of $W'_{red} \to W_{red}$. Then $V_{/Z} \times_{g, W} W' \to V_{/Z}$
is an adic \'etale morphism of formal algebraic spaces over $S$
and $V_{/Z} \times_{g, W} W'$ is an affine formal algebraic space.
By Algebraization of Formal Spaces,
Lemma \ref{restricted-lemma-algebraize-rig-etale-affine}
there exists an \'etale morphism $\varphi : V' \to V$ of affine schemes
such that the completion of $V'$ along $Z' = \varphi^{-1}(Z)$
is isomorphic to $V_{/Z} \times_{g, W} W'$ over $V_{/Z}$.
Observe that $v$ is the image of some $v' \in V'$.
Since smoothness is preserved under base change we see that
$V'_n \to W'_n$ is smooth for all $n$. In the next paragraph
we show that after replacing $V'$ by an open neighbourhood of $v'$
the morphisms $V'_n \to W'_n$ are smooth for all $n$.
Then, after we replace $V$ by the open image of $V' \to V$,
we obtain that $V_n \to W_n$ is smooth by \'etale descent of smoothness.
Some details omitted.

\medskip\noindent
Assume $S = \Spec(R)$, $V = \Spec(A)$, $Z = V(I)$, and $W = \text{Spf}(B)$.
Let $v$ correspond to the maximal ideal $I \subset \mathfrak m \subset A$.
We are given an adic continuous $R$-algebra homomorphism
$$
B \longrightarrow A^\wedge
$$
Let $\mathfrak b \subset B$ be the ideal of topologically nilpotent
elements (this is the maximal ideal of definition of the Noetherian adic
topological ring $B$). Observe that $\mathfrak b A^\wedge$ and
$IA^\wedge$ are both ideals of definition of the Noetherian adic
ring $A^\wedge$. Also, $\mathfrak m A^\wedge$ is a maximal ideal
of $A^\wedge$ containing both $\mathfrak b A^\wedge$ and $IA^\wedge$.
We are given that
$$
B_n = B/\mathfrak b^n \to A^\wedge/\mathfrak b^n A^\wedge = A_n
$$
is smooth at $\mathfrak m$ for all $n$. By the discussion above
we may and do assume that $B_1 \to A_1$ is a smooth ring map.
Denote $\mathfrak m_1 \subset A_1$ the maximal ideal corresponding
to $\mathfrak m$. Since smoothness implies flatness, we see that:
for all $n \geq 1$ the map
$$
\mathfrak b^n/\mathfrak b^{n + 1} \otimes_{B_1} (A_1)_{\mathfrak m_1}
\longrightarrow
\left(\mathfrak b^nA^\wedge/\mathfrak b^{n + 1}A^\wedge\right)_{\mathfrak m_1}
$$
is an isomorphism (see
Algebra, Lemma \href{https://stacks.math.columbia.edu/tag/0AS8}{lemma what does it mean again (Tag 0AS8)}).
Consider the Rees algebra
$$
B' = \bigoplus\nolimits_{n \geq 0} \mathfrak b^n/\mathfrak b^{n + 1}
$$
which is a finite type graded algebra over the Noetherian ring $B_1$ and
the Rees algebra
$$
A' = \bigoplus\nolimits_{n \geq 0}
\mathfrak b^nA^\wedge/\mathfrak b^{n + 1}A^\wedge
$$
which is a a finite type graded algebra over the Noetherian ring $A_1$.
Consider the homomorphism of graded $A_1$-algebras
$$
\Psi : B' \otimes_{B_1} A_1 \longrightarrow A'
$$
By the above this map is an isomorphism after localizing at
the maximal ideal $\mathfrak m_1$ of $A_1$.
Hence $\Ker(\Psi)$, resp.\ $\Coker(\Psi)$ is a finite module
over $B' \otimes_{B_1} A_1$, resp.\ $A'$ whose localization
at $\mathfrak m_1$ is zero. It follows that after replacing
$A_1$ (and correspondingly $A$) by a principal localization
we may assume $\Psi$ is an isomorphism. (This is the key step of the proof.)
Then working backwards we see that $B_n \to A_n$ is flat, see
Algebra, Lemma \href{https://stacks.math.columbia.edu/tag/0AS8}{lemma what does it mean again (Tag 0AS8)}.
Hence $A_n \to B_n$ is smooth (as a flat ring map with smooth
fibres, see Algebra, Lemma \href{https://stacks.math.columbia.edu/tag/00TF}{lemma flat fibre smooth (Tag 00TF)})
and the proof is complete.
\end{proof}
\end{document}
